% Created with jtex v.1.0.21
\documentclass{article}
\PassOptionsToPackage{short, nodayofweek}{datetime}

\input{curvenote.def}



% colors for hyperlinks
\hypersetup{colorlinks=true, allcolors=blue}
\hypersetup{
pdftitle={\@title},
pdfsubject={},
pdfauthor={\@author},
pdfkeywords={},
addtopdfcreator={Written in Curvenote}
}

\usepackage{curvenote}

\title{Hull modelling in Rhino}

\newdate{articleDate}{7}{10}{2025}
\date{\displaydate{articleDate}}

\author{\bfseries Josip Bašić\mdseries\\}

\begin{document}

\maketitle
\keywords{}

\section{Import hull lines drawing}

An example of how to import picture into Rhino is given here and explained below.

\subsection{Set up the drawing space}

\begin{enumerate}
\item Start by maximizing the desired viewport for your tracing. For example, double-click on the ``Right'' viewport title to make it the only visible viewport.
\item In the command line, type the command \textit{Picture} and press \textit{Enter}.
\item An ``\textit{Open Bitmap}'' window will appear. Navigate to your image file, select it, and click ``\textit{Open}''.
\item Click and drag in the viewport to place the image onto the construction plane.
\end{enumerate}

\subsection{Adjust image properties}

\begin{enumerate}
\item Select the image you just placed.
\item On the right-hand side, go to the \textbf{Properties} panel and select the \textbf{Material} tab (it looks like a paint tube).
\item In the ``\textit{Picture}'' settings, check the box for \textbf{Color Mask}.
\item A ``\textit{Select Color}'' window will pop up. Choose the image background color (usually it is \textbf{White)} from the color list and click \textbf{OK}. The chosen background (e.g. white colour) of your image will now be transparent.
\item If the result is a bit noisy, play with \textbf{Tolerance} (T:) values to clean out more background color.
\item To make the drawing lines easier to see against the grid, adjust the \textbf{Transparency} slider. A setting of around 70-80\% is usually effective.
\end{enumerate}

\subsection{Align the Image to the Grid:}

\begin{enumerate}
\item \textbf{Move the Image:}

\begin{itemize}
\item Type the Move command and press Enter.
\item Select a key reference point on your drawing to move \textit{from}. This should be a point you want to align with the origin (0,0,0), such as the intersection of a baseline and a center line.
\item For the point to move \textit{to}, type 0 in the command line and press \textit{Enter}. This will snap your reference point to the origin of the grid.
\end{itemize}


\item \textbf{Rotate the Image:}

\begin{itemize}
\item Type the \textit{Rotate} command and press \textit{Enter}.
\item For the center of rotation, select the same reference point you just moved to the origin.
\item For the first reference point, click on another point along the baseline of your drawing.
\item To align it perfectly horizontally, hold the \textbf{Shift} key (to enable \textit{Ortho} mode) and click anywhere along the horizontal axis. Your image's baseline will now be aligned with the grid.
\end{itemize}
\end{enumerate}

\subsection{Only if the image is scanned: undeform it}

If your scanned image is distorted or skewed, you can add control points to the image plane itself and move them to align the drawing with the grid accurately.

\begin{enumerate}
\item Select the image plane and type the \textbf{Rebuild} command and press \textbf{Enter}.
\item In the ``\textbf{Rebuild Surface}'' dialog box, you can set the \textbf{Point Count} for the U and V directions. This determines how many control points your surface will have. For example, set both U and V to \textbf{3}. For more complex adjustments, you might use a higher number. Click \textbf{OK}.
\item With the image still selected, type the \textbf{PointsOn} command to see a grid of control points.
\item Drag control points to stretch or compress that part of the image until the underlying drawing feature (like a horizontal line) aligns perfectly with the grid.
\item Repeat this process for other points as needed to correct the entire image.
\item Once you are satisfied with the alignment, press the \textbf{Esc} key to deselect the points and hide the control point grid.
\end{enumerate}

\section{Setting up modelling process}

\subsection{Scale the Drawing to the Correct Dimensions}

Before tracing, it's good to ensure the reference image is scaled correctly. If the drawing has a dimension marked on it (like the length or beam), you can use that to scale the image accurately.

\begin{enumerate}
\item Select your image plane and call the \textbf{Scale2D} command.
\item \textbf{First reference point:} Click on one end of a known dimension on your drawing (e.g., the start of the waterline).
\item \textbf{Second reference point:} Click on the other end of that same dimension.
\item You will then be prompted to enter the new distance. Type the actual length that the reference line represents and press \textbf{Enter}. Your image is now scaled correctly.
\end{enumerate}

\subsection{Organize Your Workspace with Layers}

Using layers is essential for keeping the model organized.

\begin{enumerate}
\item Open the \textbf{Layers} panel.
\item Rename the default layer to ``background''.
\item Select your image plane, right-click on the ``background'' layer, and choose \textbf{Change Object Layer}.
\item Click the lock icon next to the ``background'' layer to prevent accidental moving of the image.
\item Create a new layer and name it ``curves''. Click the circle next to it to make it the active layer. All new geometry will now be created on this layer.
\end{enumerate}

\section{Drawing sections}

An example how to draw sections and produce hull surface is given in the video here and explained below.

\subsection{Trace the hull sections}

Be sure to maximize the window showing the hull lines. Now each section curve needs to be digitally drawn by using commands from the \textbf{Curves} toolbar.

\begin{enumerate}
\item From the toolbar, select the \textbf{Curve: Interpolate Points} tool.
\item Carefully click along the contour of one of the station lines on your drawing. Place points at key locations where the curve changes direction.
\item Right-click or press Enter to complete the curve.
\item Repeat this process for all the visible station lines on one side of the hull.
\end{enumerate}

Now let's ensure all endpoints of drawn curves start for the middle of the ship (i.e. centreline).

\begin{enumerate}
\item Select all or one-by-one drawn section curves.
\item Use the \textbf{PointsOn} command (or \textbf{F10}) to show its control points.
\item Select all or one-by-one control points nearest to the centreline.
\item Type \textbf{SetPt} (Set Points) and press \textbf{Enter}. In the dialog box, check \textbf{Set Y} and \textbf{Align to World}, then click OK. Enter 0 for the \textbf{Location for points}. Now all curves start at \textit{Y}=0, i.e. the centreline.
\item Alternatively to the step above, you can enable \textbf{Grid Snap} and \textbf{Ortho} to simply drag the control points to the centreline.
\end{enumerate}

\subsection{Align the curves along longitudinal axis}

The curves you drew are 2D, but they represent sections at different positions along the hull. We need to move them into their correct 3D locations.

\begin{enumerate}
\item Switch to the \textbf{Front} viewport. You will see all your traced curves stacked on top of each other.
\item By using the \textbf{Gumball} or \textbf{Move} command, place each curve to its designated station location along the longitudinal axis.
\item Use the \textbf{Top} and \textbf{Front} viewports to review the process and to guide you.
\end{enumerate}

\subsection{Create the keel and chine curves (rails)}

To create the keel curve, we have to connect all lowest ends of section curves.

\begin{enumerate}
\item In the \textbf{Right} viewport, enable \textbf{OSnap}, and call the \textbf{Curve: Interpolate Points} tool again.
\item Create a curve from the first to last section, by clicking on the snapped lowest points of each section curve.
\item Right-click or press Enter to complete the curve, when all keel points are connected.
\end{enumerate}

If the ship contains a chine, repeat the above process to create the chine curve.
\end{document}
